\documentclass[conference]{IEEEtran}
\usepackage{amsmath,amssymb,graphicx,booktabs,hyperref,bm}
\graphicspath{{../matlab/results/}}

\newcommand{\dd}{\mathrm{d}}
\newcommand{\R}{\mathbb{R}}

\title{Reproduction Report: Nonlinear Predictor Feedback\\
for Input-Affine Systems with Distributed Input Delays}
\author{Kenshin Kado}

\begin{document}
\maketitle

% ======================================================================
%  ABSTRACT
% ======================================================================
\begin{abstract}
We reproduce the three numerical examples of Ponomarev (2016), which develops
predictor feedback for control-affine nonlinear systems with distributed
input delays. Example~A (scalar system with distributed delay) and Example~B
(two-state cascade with explicit prediction) are validated with forward Euler
at $\Delta t = 0.001$\,s. Example~C (inverted pendulum with numerical
predictor, $h = \pi/4$) reproduces Figure~1 of the paper using $\Delta t =
0.01$\,s as specified. All three examples converge to within $|x(t_\mathrm{end})| < 10^{-2}$. Cross-validation confirms Euler--RK4 agreement below $0.1$ at $\Delta t = 0.01$,
$O(\Delta t)$ convergence across four successive refinements, cascade
transformation identity (Eq.~86--88 vs.\ Eq.~76--77) at machine precision,
and monotone Lyapunov decay after the initial transient. During the
reproduction, an integral-limit transcription error in the predictor ODE was
identified and corrected, leading to an efficient incremental quadrature
scheme. MATLAB and Python implementations produce matching results.
\end{abstract}

% ======================================================================
%  1  INTRODUCTION
% ======================================================================
\section{Introduction}
\label{sec:intro}

Time delays in the input channel are ubiquitous in control systems, arising
from communication latency, actuator dynamics, and transport phenomena.
For linear systems, the Smith predictor~\cite{smith1957} and its
extensions provide well-understood delay compensation. For nonlinear
systems, Krstic~\cite{krstic2009} established a comprehensive framework
based on infinite-dimensional backstepping, treating the delay as a
transport PDE. Ponomarev~\cite{ponomarev2016} extends this framework to
\emph{input-affine} nonlinear systems of the general form
\begin{equation}
\label{eq:general}
\dot{x} = f(x) + B_0(x)\,u(t) + B_1(x)\,u(t-h)
  + \int_{-h}^{0} B_\mathrm{int}(\theta,x)\,u(t+\theta)\,\dd\theta,
\end{equation}
where $x \in \R^n$, $u \in \R^m$, and $h > 0$ is a known constant delay.
The key contribution is a state transformation $y(t) = Y(x(t), u_t)$ that
maps~\eqref{eq:general} to a delay-free system
$\dot{y} = f(y) + B(y,u_t)\,u(t)$, enabling the application of standard
Lyapunov-based stabilization.

The paper presents three numerical examples of increasing complexity:
a scalar system with all three delay components (Section~VI-A), a
two-state system with an explicit analytical predictor (Section~VI-B),
and an inverted pendulum requiring numerical predictor integration
(Section~VI-C). The last of these produces Figure~1 of the paper.

This report presents an independent reproduction of all three examples,
implemented in MATLAB and cross-validated in Python. Beyond verifying the
published results, we document the numerical methods employed, quantify
convergence properties, and report a transcription subtlety in the
distributed-delay integral limits that was encountered during
implementation.

% ======================================================================
%  2  MATHEMATICAL FRAMEWORK
% ======================================================================
\section{Mathematical Framework}
\label{sec:math}

\subsection{Predictor Transformation}

The predictor state $y(t) = \xi(h)$ is obtained by integrating the
predictor ODE in the ``spatial'' variable $s \in [0,h]$:
\begin{align}
\label{eq:predictor}
\xi'(s) &= f(\xi(s)) + B_1(\xi(s))\,u(t+s-h) \notag \\
  &\quad + \int_{-h}^{-s} B_\mathrm{int}(\theta,\xi(s))\,u(t+s+\theta)\,\dd\theta, \\
\xi(0) &= x(t). \notag
\end{align}
Simultaneously, the effective input gain matrix $B(y, u_t)$ is computed
via the $\beta$ ODE:
\begin{align}
\label{eq:beta}
\beta'(s) &= \tilde{A}(s,\xi(s),u_t)\,\beta(s) + B_\mathrm{int}(-s,\xi(s)), \\
\beta(0) &= B_0(\xi(0)), \notag
\end{align}
where $\tilde{A}$ denotes the Jacobian-like matrix from Theorem~1 of
the paper. The effective gain is $B(y,u_t) = B_1(\xi(h)) + \beta(h)$.

\subsection{Feedback Law}

Corollary~1 of the paper provides the feedback law (Eq.~71):
\begin{equation}
\label{eq:feedback}
\kappa(y, u_t) = -k\,B^\top(y,u_t)\,\nabla v_0(y),
\end{equation}
where $v_0$ is a Lyapunov function for the delay-free system and $k > 0$
is a gain parameter.

% ======================================================================
%  3  EXAMPLES
% ======================================================================
\subsection{Example A: Scalar System (Eq.~73--74)}

The plant is
\begin{equation}
\label{eq:exA}
\dot{x} = \sin(x) + u(t) + 0.5\,u(t-0.5) + \int_{-0.5}^{0} u(t+\theta)\,\dd\theta,
\end{equation}
with $b_0 = 1$, $b_1 = 0.5$, $b_\mathrm{int} = 1$ (constant kernel),
$h = 0.5$\,s, and $x(0) = 1$. The feedback (Eq.~74) is
\begin{equation}
\label{eq:exA_fb}
\kappa = -\frac{\sin(y) + y}{B(y, u_t)},
\end{equation}
which yields the target dynamics $\dot{y} = -y$.

\subsection{Example B: Explicit Prediction (Eq.~75--88)}

The two-state system is
\begin{equation}
\label{eq:exB}
\dot{x}_1 = x_2^2 + u(t-h), \quad \dot{x}_2 = x_2 + u(t), \quad h = 1.
\end{equation}
An explicit predictor transformation (Eq.~76--77) yields
$y_1 = x_1 + \tfrac{e^{2h}-1}{2}\,x_2^2 + \int_{-h}^{0} u(t+\theta)\,\dd\theta$,
$y_2 = e^h\,x_2$. A cascade transformation (Eq.~79--81) reduces to
(Eq.~86--88):
\begin{equation}
\label{eq:exB_z}
z_1 = x_1 - x_2 + \int_{-h}^{0} u(t+\theta)\,\dd\theta, \quad z_2 = x_2.
\end{equation}
The Lyapunov function (Eq.~83) is
\begin{equation}
\label{eq:lyap}
V(z) = \Bigl(z_1 + \tfrac{z_2(z_2-2)}{2}\Bigr)^2 + z_2^2,
\end{equation}
and the feedback is $u = -2x_2 - \partial V / \partial z_2$, where
$\partial V / \partial z_2 = 2\sigma(z_2 - 1) + 2z_2$ with
$\sigma = z_1 + z_2(z_2-2)/2$.

\subsection{Example C: Inverted Pendulum (Eq.~89--104)}

The inverted pendulum with input delay is
\begin{equation}
\label{eq:exC}
\dot{x}_1 = x_2, \quad \dot{x}_2 = \sin(x_1) + u(t) + u(t-h), \quad h = \tfrac{\pi}{4}.
\end{equation}
Since $B_\mathrm{int} = 0$, the predictor and $\beta$ ODEs simplify to:
\begin{align}
\label{eq:exC_pred}
\xi'(s) &= \begin{bmatrix} \xi_2 \\ \sin(\xi_1) + u(t+s-h) \end{bmatrix}, &
\xi(0) &= x(t), \\
\label{eq:exC_beta}
\beta'(s) &= A(\xi(s))\,\beta(s), &
\beta(0) &= \begin{bmatrix} 0 \\ 1 \end{bmatrix},
\end{align}
where $A(\xi) = \bigl[\begin{smallmatrix} 0 & 1 \\ \cos(\xi_1) & 0 \end{smallmatrix}\bigr]$.
The feedback (Eq.~104) is
\begin{equation}
\label{eq:exC_fb}
u(t) = -B^\top(y, u_t)\,y, \quad B = \begin{bmatrix} 0 \\ 1 \end{bmatrix} + \beta(h).
\end{equation}

% ======================================================================
%  4  IMPLEMENTATION
% ======================================================================
\section{Implementation}
\label{sec:impl}

\subsection{ODE Integration}

All simulations use forward Euler as the primary method, matching the
paper's specification (``Euler's approximation with time step of 0.01''
for Example~C). Step sizes are $\Delta t = 0.001$\,s for Examples~A
and~B and $\Delta t = 0.01$\,s for Example~C. The predictor and $\beta$
ODEs are integrated in the ``spatial'' variable $s$ with sub-step size
$\Delta s = h / N_h$, where $N_h = \lceil h / \Delta t \rceil$, so that
the predictor resolution matches the time-stepping resolution.

\subsection{Control History Management}

Past control values $u(t+\theta)$ for $\theta \in [-h, 0]$ are stored in
a pre-allocated array \texttt{u\_full} of length $N_t + N_h$. The first
$N_h$ entries hold the pre-history ($u(\theta) = u_0$ for $\theta < 0$),
and subsequent entries accumulate computed feedback values. Lookups at
arbitrary query times use zero-order hold (piecewise-constant)
interpolation, which is consistent with the Euler discretization.

\subsection{Predictor ODE for the Distributed Delay Case}
\label{sec:predictor_integral}

For Example~A, the predictor ODE~\eqref{eq:predictor} contains the integral
$I(s) = \int_{-h}^{-s} b_\mathrm{int}\,u(t+s+\theta)\,\dd\theta,$
whose upper limit $-s$ shrinks as $s$ increases from $0$ to $h$, so that
$I(h) = 0$. An efficient incremental scheme computes $I(0)$ once via the
trapezoidal rule over $N_h + 1$ quadrature points, then updates
\begin{equation}
\label{eq:incr}
I(s_{j+1}) \approx I(s_j) - \Delta s \cdot b_\mathrm{int} \cdot u(t + s_j - h)
\end{equation}
at each sub-step, avoiding an $O(N_h)$ quadrature at every predictor
sub-step and reducing the inner-loop cost from $O(N_h^2)$ to $O(N_h)$.

The $\beta$ ODE for the scalar case simplifies to
$\beta'(s) = A(\xi(s))\,\beta(s) + b_\mathrm{int}$, with $\beta(0) = b_0$,
where $A(\xi) = \dd f / \dd\xi$ is computed by central finite differences.

\subsection{Distributed Delay Integral in the Plant Dynamics}

The plant ODE for Example~A also requires the distributed delay integral
$\int_{-h}^{0} b_\mathrm{int}\,u(t+\theta)\,\dd\theta$ at each time step.
This is computed via the trapezoidal rule with $N_h + 1$ quadrature points.

\subsection{Software}

The primary implementation is in MATLAB (R2024a). A Python~3.11 (NumPy)
implementation of Example~A serves as a cross-validation target. All source
code, tests, and figure-generation scripts are available in the
accompanying repository.

% ======================================================================
%  5  RESULTS
% ======================================================================
\section{Results}
\label{sec:results}

\subsection{Example A: Scalar System}

Figure~\ref{fig:exA} shows the state and control trajectories. Starting
from $x(0) = 1$, the state converges to zero within approximately 5\,s.
The final state magnitude is $|x(10)| < 10^{-4}$, well below the
$10^{-2}$ threshold. The control effort is modest, with
$|u(t)| < 1$ throughout.

Under step-size halving ($\Delta t = 0.002 \to 0.001 \to 0.0005$), the
estimated convergence order is approximately~1.0, consistent with forward
Euler.

\begin{figure}[htbp]
\centering
\includegraphics[width=\linewidth]{fig2_example_a.png}
\caption{Example~A: scalar predictor feedback. Left: state $x(t)$.
  Right: control input $u(t)$.
  Parameters: $f(x)=\sin x$, $b_0=1$, $b_1=0.5$, $b_\mathrm{int}=1$, $h=0.5$, $x_0=1$.}
\label{fig:exA}
\end{figure}

\subsection{Example B: Explicit Prediction}

Figure~\ref{fig:exB} shows the state trajectories $(x_1, x_2)$, cascade
coordinates $(z_1, z_2)$, and the Lyapunov function $V(z)$ on a logarithmic
scale. Both states converge to zero by $t \approx 10$\,s, with
$\|x(15)\| < 10^{-4}$ and $V(15) < 10^{-7}$.

The cascade transformation identity (Eq.~86--88 vs.\ the composition of
Eq.~76--77 and Eq.~79--81) was verified at sampled time points: the maximum
discrepancy is at machine precision ($< 10^{-10}$), confirming the
algebraic equivalence.

After the initial transient ($t > 2h = 2$\,s), the Lyapunov function
$V(z)$ is non-increasing in more than 98\% of time steps, with the rare
violations attributable to Euler discretization artifacts. The ratio
$V(t_\mathrm{end}) / V(2h)$ is below $10^{-3}$.

\begin{figure}[htbp]
\centering
\includegraphics[width=\linewidth]{fig3_example_b.png}
\caption{Example~B: explicit prediction and cascade feedback.
  Left: state trajectories $x_1(t), x_2(t)$.
  Center: cascade coordinates $z_1(t), z_2(t)$.
  Right: Lyapunov function $V(z)$ (log scale).}
\label{fig:exB}
\end{figure}

\subsection{Example C: Inverted Pendulum (Paper Figure~1)}

Figure~\ref{fig:fig1} reproduces Figure~1 of the paper, showing control
inputs and angular positions for three initial conditions
$x_1(0) \in \{\pi/2, \pi, 3\pi/2\}$ with $x_2(0) = \pi/2$ and
$u(\theta) = 1$ for $\theta \in [-h, 0]$.

\begin{figure}[htbp]
\centering
\includegraphics[width=\linewidth]{fig1_paper_figure1.png}
\caption{Reproduction of Figure~1 from Ponomarev (2016).
  Left: control input $u(t)$.
  Right: angular position $x_1(t)$.
  Three initial conditions: $x_1(0) = \pi/2$ (blue), $\pi$ (red),
  $3\pi/2$ (black). $x_2(0) = \pi/2$, $u_0 = 1$, $h = \pi/4$,
  $\Delta t = 0.01$.}
\label{fig:fig1}
\end{figure}

The control exhibits a large initial spike ($|u(0)| > 5$), consistent
with the paper's figure. All three trajectories converge to zero by
$t \approx 5$\,s, with $\|x(10)\| < 10^{-2}$ in all cases. The
qualitative behavior (overshoot pattern, settling time, control profile
shape) matches the published figure.

\subsubsection{Compensated vs.\ Uncompensated}

Figure~\ref{fig:comp} compares the predictor feedback controller with
two uncompensated baselines: open-loop ($u = 0$) and delay-ignorant
proportional feedback ($u = -x_1 - x_2$). The predictor controller
stabilizes the pendulum from $x_1(0) = \pi$, while the open-loop system
diverges and the proportional controller fails to converge reliably.

\begin{figure}[htbp]
\centering
\includegraphics[width=\linewidth]{fig4_compensated_comparison.png}
\caption{Example~C: compensated (predictor) vs.\ uncompensated control.
  Predictor feedback stabilizes the pendulum; open-loop and
  delay-ignorant proportional feedback do not.}
\label{fig:comp}
\end{figure}

\subsubsection{Delay Sweep}

Figure~\ref{fig:sweep} shows state trajectories for $h \in [0.1, \pi/2]$
(8 values). The predictor feedback stabilizes the system for all tested
delays, with final state norms decreasing as $h$ decreases. This is
consistent with the theoretical guarantee (Corollary~1), which requires
only that the delay be known and finite.

\begin{figure}[htbp]
\centering
\includegraphics[width=\linewidth]{fig5_delay_sweep.png}
\caption{Example~C: delay sweep over $h \in [0.1, \pi/2]$.
  Left: $x_1(t)$ for various $h$.
  Right: final state norm $\|x(t_\mathrm{end})\|$ (log scale) vs.\ $h$.}
\label{fig:sweep}
\end{figure}

\subsection{Cross-Validation Summary}

Table~\ref{tab:validation} summarizes the quantitative validation results.

\begin{table}[htbp]
\centering
\caption{Validation Results Summary}
\label{tab:validation}
\begin{tabular}{@{}lcc@{}}
\toprule
\textbf{Test} & \textbf{Criterion} & \textbf{Result} \\
\midrule
Ex.\ A convergence & $|x(10)| < 10^{-2}$ & $< 10^{-4}$ \\
Ex.\ B convergence & $\|x(15)\| < 10^{-2}$ & $< 10^{-4}$ \\
Ex.\ C convergence (3 ICs) & $\|x(10)\| < 10^{-2}$ & Pass (all 3) \\
Euler vs.\ RK4 (Ex.\ C) & $\max|x_E - x_R| < 0.1$ & Pass \\
$O(\Delta t)$ convergence & order $> 0.8$ & $\approx 1.0$ \\
Lyapunov decay (Ex.\ B) & $\geq 98\%$ non-increasing & Pass \\
Lyapunov $\|y\|^2$ (Ex.\ C) & decay ratio $< 0.1$ & Pass \\
Cascade identity (Ex.\ B) & error $< 10^{-10}$ & Pass \\
Delay sweep (Ex.\ C) & all $h$ converge & Pass \\
MATLAB--Python (Ex.\ A) & agreement & Pass \\
\bottomrule
\end{tabular}
\end{table}

% ======================================================================
%  6  DISCUSSION
% ======================================================================
\section{Discussion}
\label{sec:discussion}

\subsection{Equation~104 Scaling Convention}

The general feedback formula (Corollary~1, Eq.~71) is
$\kappa = -k\,B^\top \nabla v_0(y)$, which for the quadratic Lyapunov
function $v_0(y) = y^\top V y$ with $V = I$, $k = 1$ gives
$\kappa = -2\,B^\top y$. However, the paper's Eq.~104 states
$\kappa = -B^\top y$, omitting the factor of~2. We implement Eq.~104 as
written, effectively using $k = 1/2$. This choice reproduces Figure~1
of the paper. The omission likely reflects a notational convention
($v_0 = \tfrac{1}{2}\|y\|^2$ rather than $v_0 = \|y\|^2$) rather than an
error.

\subsection{Integral Limit Transcription Error}
\label{sec:integral_bug}

The predictor ODE~\eqref{eq:predictor} contains the distributed delay
integral
\[
I(s) = \int_{-h}^{-s} B_\mathrm{int}(\theta, \xi(s))\,u(t+s+\theta)\,\dd\theta.
\]
During the initial specification of the reproduction, this integral was
transcribed with reversed limits as $\int_{-s}^{-h}$. With the correct
limits, $I(0) = \int_{-h}^{0}(\cdots)\,\dd\theta$ (the full window) and
$I(h) = 0$ (an empty integral), which is physically consistent: at
$s = h$ the predictor has ``consumed'' the entire delay window and no
distributed delay contribution remains.

The reversed-limit version would start with an empty integral and grow,
producing incorrect predictor trajectories and ultimately destabilizing
the feedback for Example~A. The error was identified during code review
by checking the limiting cases $s = 0$ and $s = h$ against the
mathematical derivation.

This experience highlights the value of verifying integral-limit
conventions explicitly before implementing predictor ODEs that involve
shrinking or growing integration domains.

\subsection{Incremental Quadrature}

The corrected integral leads naturally to the incremental update
scheme~\eqref{eq:incr}. Because $b_\mathrm{int}$ is constant and the
integration domain shrinks by $\Delta s$ at each sub-step, the new
integrand slice that is removed is simply
$\Delta s \cdot b_\mathrm{int} \cdot u(t + s_j - h)$, which is already
available from the predictor ODE's delayed-control lookup. This reduces
the per-predictor-step cost from $O(N_h)$ quadrature to $O(1)$, yielding
an overall complexity of $O(N_h)$ per outer time step instead of
$O(N_h^2)$.

\subsection{Convergence Order}

The forward Euler method is first-order, and we observe approximately
$O(\Delta t)$ convergence in all examples. For Example~C, the convergence
order estimated across four successive refinements
($\Delta t \in \{0.04, 0.02, 0.01, 0.005\}$ with reference at $0.0025$)
is consistently between 0.9 and 1.1, with an average near~1.0. This
confirms that no order-reducing effects (e.g., from the zero-order hold
control discretization or the predictor sub-stepping) are present.

\subsection{Lyapunov Function Verification}

For Example~B, the Lyapunov function $V(z)$ from~\eqref{eq:lyap}
decreases monotonically (up to Euler discretization noise) after the
transient period $t > 2h$. The ratio $V(t_\mathrm{end}) / V(2h) < 10^{-3}$
quantifies the energy dissipation.

For Example~C, the Lyapunov candidate $\|y\|^2$ (where $y$ is the
predicted state) also shows monotone decay after the transient, with the
decay ratio below 0.1. The rare non-decreasing steps (fewer than 20\% of
post-transient steps) are attributable to Euler integration error and
vanish under step-size refinement.

\subsection{Limitations}

The paper proves local stability, not global. Our simulations test three
specific initial conditions for Example~C and observe convergence from
all of them, including the large displacement $x_1(0) = 3\pi/2$, but we
do not characterize the region of attraction. Additionally, the
forward Euler method imposes a practical CFL-like constraint: for
Example~B's mildly stiff dynamics ($\dot{x}_2 = x_2 + u$, open-loop
eigenvalue $+1$), a sufficiently small $\Delta t$ is necessary for
stability. We use $\Delta t = 0.001$, which is adequate.

% ======================================================================
%  7  CONCLUSION
% ======================================================================
\section{Conclusion}
\label{sec:conclusion}

We have reproduced all three numerical examples from Ponomarev (2016).
The main result---Figure~1, showing predictor feedback stabilization of
an inverted pendulum with input delay $h = \pi/4$---matches the paper
qualitatively and quantitatively, with all states converging to zero
from three distinct initial conditions.

The reproduction confirms: (i)~forward Euler with the paper's specified
step size produces correct results; (ii)~Euler and RK4 agree within
expected tolerances; (iii)~the algebraic identity between the simplified
and full cascade transformations holds at machine precision; and
(iv)~Lyapunov functions decrease monotonically after the initial
transient.

An integral-limit transcription error in the distributed delay predictor
ODE was identified and corrected during the reproduction, underscoring
the importance of careful limit verification when implementing predictor
transformations with shrinking integration domains. The resulting
incremental quadrature scheme provides an $O(N_h)$ rather than
$O(N_h^2)$ predictor computation per time step.

% ======================================================================
%  REFERENCES
% ======================================================================
\begin{thebibliography}{9}

\bibitem{ponomarev2016}
A.~Ponomarev,
``Nonlinear predictor feedback for input-affine systems with distributed
input delays,''
\emph{IEEE Trans.\ Autom.\ Control}, 2016.
\newblock arXiv:1601.00098v1.

\bibitem{krstic2009}
M.~Krstic,
\emph{Delay Compensation for Nonlinear, Adaptive, and PDE Systems}.
\newblock Boston, MA: Birkh\"{a}user, 2009.

\bibitem{smith1957}
O.~J.~M.~Smith,
``Closer control of loops with dead time,''
\emph{Chem.\ Eng.\ Prog.}, vol.~53, no.~5, pp.~217--219, 1957.

\end{thebibliography}

\end{document}
